% GENERATED FILE - do not hand-edit.
% Source of truth: data/calculus-reviewed/t-1.json in the zenith-website repo.
% Regenerate with: python3 scripts/build_calculus_tex.py
\documentclass[11pt]{article}
\usepackage[utf8]{inputenc}
\usepackage[T1]{fontenc}
\usepackage{amsmath,amssymb}
\usepackage{graphicx}
\usepackage[margin=1in]{geometry}
\graphicspath{{../..//}}
\setlength{\parindent}{0pt}
\setlength{\parskip}{0.6em}
\begin{document}
\section*{Chapter 1: Functions --- Test T1}

\bigskip\noindent\textbf{T1.}\ If $f(x) = x^3 - 3x - 1$, then $f(-2) =$
\begin{list}{}{\setlength{\leftmargin}{3em}\setlength{\labelwidth}{2.5em}\setlength{\itemsep}{0pt}\setlength{\topsep}{2pt}}
\item[(A)] $-13$
\item[(B)] $-3$
\item[(C)] $-1$
\item[(D)] $7$
\end{list}

\bigskip\noindent\textbf{T2.}\ The domain of $f(x) = \dfrac{x+2}{x^2+4}$ is
\begin{list}{}{\setlength{\leftmargin}{3em}\setlength{\labelwidth}{2.5em}\setlength{\itemsep}{0pt}\setlength{\topsep}{2pt}}
\item[(A)] all $x \neq 2$
\item[(B)] all $x \neq 2, -2$
\item[(C)] all $x \neq -2$
\item[(D)] all reals
\end{list}

\bigskip\noindent\textbf{T3.}\ The domain of $g(x) = \dfrac{\sqrt{x-5}}{x^2-9}$ is
\begin{list}{}{\setlength{\leftmargin}{3em}\setlength{\labelwidth}{2.5em}\setlength{\itemsep}{0pt}\setlength{\topsep}{2pt}}
\item[(A)] all $x \neq 3, -3$
\item[(B)] $x \leq 5, x \neq 3, -3$
\item[(C)] $x \geq 5$
\item[(D)] $x > 5$
\end{list}

\bigskip\noindent\textbf{T4.}\ If $f(x) = x^3 - 2x^2 - 3x + 4$ and $g(x) = 5$, then $g(f(x)) =$
\begin{list}{}{\setlength{\leftmargin}{3em}\setlength{\labelwidth}{2.5em}\setlength{\itemsep}{0pt}\setlength{\topsep}{2pt}}
\item[(A)] $5x^3 - 10x^2 - 15x + 20$
\item[(B)] $5x^2 - 10x - 1$
\item[(C)] $64$
\item[(D)] $5$
\end{list}

\bigskip\noindent\textbf{T5.}\ If $f(x) = x^3 - 2x^2 - 3x + 4$ and $g(x) = 5$, then $f(g(x)) =$
\begin{list}{}{\setlength{\leftmargin}{3em}\setlength{\labelwidth}{2.5em}\setlength{\itemsep}{0pt}\setlength{\topsep}{2pt}}
\item[(A)] $5x^3 - 10x^2 - 15x + 20$
\item[(B)] $5x^2 - 10x - 1$
\item[(C)] $64$
\item[(D)] $5$
\end{list}

\bigskip\noindent\textbf{T6.}\ If $f(x) = x^3 + Ax^2 + Bx + 2$ and if $f(1) = 6$ and $f(-1) = -2$, what is the value of $2A + B$?
\begin{list}{}{\setlength{\leftmargin}{3em}\setlength{\labelwidth}{2.5em}\setlength{\itemsep}{0pt}\setlength{\topsep}{2pt}}
\item[(A)] $12$
\item[(B)] $3$
\item[(C)] $0$
\item[(D)] $-2$
\end{list}

\bigskip\noindent\textbf{T7.}\ Which of the following equations has a graph that is symmetric with respect to the origin?
\begin{list}{}{\setlength{\leftmargin}{3em}\setlength{\labelwidth}{2.5em}\setlength{\itemsep}{0pt}\setlength{\topsep}{2pt}}
\item[(A)] $y = \dfrac{x+1}{x}$
\item[(B)] $y = 3x^4 - 2$
\item[(C)] $y = x^3 + 3x$
\item[(D)] $y = x^3 - 4$
\end{list}

\bigskip\noindent\textbf{T8.}\ Let $g$ be a function defined for all reals. Which of the following conditions is not sufficient to guarantee that $g$ has an inverse function?
\begin{list}{}{\setlength{\leftmargin}{3em}\setlength{\labelwidth}{2.5em}\setlength{\itemsep}{0pt}\setlength{\topsep}{2pt}}
\item[(A)] $g(x) = mx + k,\ m \neq 0$
\item[(B)] $g$ is strictly increasing
\item[(C)] $g$ is an odd function
\item[(D)] $g$ is one-to-one
\end{list}

\bigskip\noindent\textbf{T9.}\ Let $y = f(x) = \cos(\arctan x)$. Then the range of $f$ is
\begin{list}{}{\setlength{\leftmargin}{3em}\setlength{\labelwidth}{2.5em}\setlength{\itemsep}{0pt}\setlength{\topsep}{2pt}}
\item[(A)] $\{y \mid 0 < y < 1\}$
\item[(B)] $\{y \mid 0 < y \leq 1\}$
\item[(C)] $\{y \mid -\tfrac{\pi}{2} < y < \tfrac{\pi}{2}\}$
\item[(D)] $\{y \mid -\tfrac{\pi}{2} \leq y \leq \tfrac{\pi}{2}\}$
\end{list}

\bigskip\noindent\textbf{T10.}\ Let $g(x) = |\sin x - 1|$. The maximum value attained by $g$ on the closed interval $[0, 2\pi]$ is for $x$ equal to
\begin{list}{}{\setlength{\leftmargin}{3em}\setlength{\labelwidth}{2.5em}\setlength{\itemsep}{0pt}\setlength{\topsep}{2pt}}
\item[(A)] $0$
\item[(B)] $\dfrac{\pi}{2}$
\item[(C)] $\pi$
\item[(D)] $\dfrac{3\pi}{2}$
\end{list}

\bigskip\noindent\textbf{T11.}\ Which of the following functions is not odd?
\begin{list}{}{\setlength{\leftmargin}{3em}\setlength{\labelwidth}{2.5em}\setlength{\itemsep}{0pt}\setlength{\topsep}{2pt}}
\item[(A)] $f(x) = \sin 3x$
\item[(B)] $f(x) = \tan x$
\item[(C)] $f(x) = x^3 - 2$
\item[(D)] $f(x) = \dfrac{x}{x^2+4}$
\end{list}

\bigskip\noindent\textbf{T12.}\ If the solutions to the equation $f(x) = 0$ are $x = 3, -4$, then $f(3x) = 0$ at $x =$
\begin{list}{}{\setlength{\leftmargin}{3em}\setlength{\labelwidth}{2.5em}\setlength{\itemsep}{0pt}\setlength{\topsep}{2pt}}
\item[(A)] $1$ and $-\dfrac{4}{3}$
\item[(B)] $-1$ and $\dfrac{4}{3}$
\item[(C)] $9$ and $-12$
\item[(D)] $-9$ and $12$
\end{list}

\bigskip\noindent\textbf{T13.}\ The set of zeros of $f(x) = x^3 - 6x^2 + 9x$ is
\begin{list}{}{\setlength{\leftmargin}{3em}\setlength{\labelwidth}{2.5em}\setlength{\itemsep}{0pt}\setlength{\topsep}{2pt}}
\item[(A)] $\{3\}$
\item[(B)] $\{0,3\}$
\item[(C)] $\{0,-3\}$
\item[(D)] $\{-3\}$
\end{list}

\bigskip\noindent\textbf{T14.}\ The values of $x$ for which the graphs of $y = x + 3$ and $y^2 = 4x$ intersect are
\begin{list}{}{\setlength{\leftmargin}{3em}\setlength{\labelwidth}{2.5em}\setlength{\itemsep}{0pt}\setlength{\topsep}{2pt}}
\item[(A)] $-3$ and $3$
\item[(B)] $-3$
\item[(C)] $3$
\item[(D)] no intersection
\end{list}

\bigskip\noindent\textbf{T15.}\ The function whose graph is a reflection in the $y$-axis of the graph of $f(x) = 2 - 5^x$ is
\begin{list}{}{\setlength{\leftmargin}{3em}\setlength{\labelwidth}{2.5em}\setlength{\itemsep}{0pt}\setlength{\topsep}{2pt}}
\item[(A)] $g(x) = 2 - 5^{-x}$
\item[(B)] $g(x) = 5^x - 2$
\item[(C)] $g(x) = \log_5(x-2)$
\item[(D)] $g(x) = \log_5(2-x)$
\end{list}

\bigskip\noindent\textbf{T16.}\ Let $g(x)$ be the inverse function of $f(x)$. Then $g(f(x)) =$
\begin{list}{}{\setlength{\leftmargin}{3em}\setlength{\labelwidth}{2.5em}\setlength{\itemsep}{0pt}\setlength{\topsep}{2pt}}
\item[(A)] $1$
\item[(B)] $x$
\item[(C)] $\dfrac{1}{x}$
\item[(D)] $f(x) \cdot g(x)$
\end{list}

\bigskip\noindent\textbf{T17.}\ The function $f(x) = 3x^3 - 2x - 6$ has exactly one real zero. It is between
\begin{list}{}{\setlength{\leftmargin}{3em}\setlength{\labelwidth}{2.5em}\setlength{\itemsep}{0pt}\setlength{\topsep}{2pt}}
\item[(A)] $-1$ and $0$
\item[(B)] $0$ and $1$
\item[(C)] $1$ and $2$
\item[(D)] $2$ and $3$
\end{list}

\bigskip\noindent\textbf{T18.}\ The period of $f(x) = \sin\!\left(\dfrac{3\pi}{4}x\right)$ is
\begin{list}{}{\setlength{\leftmargin}{3em}\setlength{\labelwidth}{2.5em}\setlength{\itemsep}{0pt}\setlength{\topsep}{2pt}}
\item[(A)] $\dfrac{3}{8}$
\item[(B)] $\dfrac{4}{3}$
\item[(C)] $\dfrac{8}{3}$
\item[(D)] $\dfrac{16}{3}$
\end{list}

\bigskip\noindent\textbf{T19.}\ The range of $y = f(x) = \ln(\sin x)$, for $0 < x < \pi$, is
\begin{list}{}{\setlength{\leftmargin}{3em}\setlength{\labelwidth}{2.5em}\setlength{\itemsep}{0pt}\setlength{\topsep}{2pt}}
\item[(A)] $\{y \mid -\infty < y \leq 0\}$
\item[(B)] $\{y \mid 0 < y \leq 1\}$
\item[(C)] $\{y \mid -1 < y < 1\}$
\item[(D)] $\{y \mid 0 \leq y \leq 1\}$
\end{list}

\bigskip\noindent\textbf{T20.}\ If $\log_b(4^b) = \dfrac{b}{2}$, then $b =$
\begin{list}{}{\setlength{\leftmargin}{3em}\setlength{\labelwidth}{2.5em}\setlength{\itemsep}{0pt}\setlength{\topsep}{2pt}}
\item[(A)] $\dfrac{1}{16}$
\item[(B)] $\dfrac{1}{4}$
\item[(C)] $4$
\item[(D)] $16$
\end{list}

\bigskip\noindent\textbf{T21.}\ Let $f^{-1}$ be the inverse function of $f(x) = x^3 - 5$. Then $f^{-1}(x) =$
\begin{list}{}{\setlength{\leftmargin}{3em}\setlength{\labelwidth}{2.5em}\setlength{\itemsep}{0pt}\setlength{\topsep}{2pt}}
\item[(A)] $\dfrac{1}{x^3+5}$
\item[(B)] $(x+5)^3$
\item[(C)] $\sqrt[3]{x-5}$
\item[(D)] $\sqrt[3]{x+5}$
\end{list}

\bigskip\noindent\textbf{T22.}\ The set of $x$-intercepts of the graph of $f(x) = x^3 - 3x^2 - x + 3$ is
\begin{list}{}{\setlength{\leftmargin}{3em}\setlength{\labelwidth}{2.5em}\setlength{\itemsep}{0pt}\setlength{\topsep}{2pt}}
\item[(A)] $\{-1,1\}$
\item[(B)] $\{1,3\}$
\item[(C)] $\{-1,1,3\}$
\item[(D)] $\{-1,-3,3\}$
\end{list}

\bigskip\noindent\textbf{T23.}\ If the domain of $f$ is restricted to the open interval $\left(-\tfrac{\pi}{2}, \tfrac{\pi}{2}\right)$, then the range of $f(x) = 5^{\tan x}$ is
\begin{list}{}{\setlength{\leftmargin}{3em}\setlength{\labelwidth}{2.5em}\setlength{\itemsep}{0pt}\setlength{\topsep}{2pt}}
\item[(A)] the set of all reals
\item[(B)] the set of positive reals
\item[(C)] the set of nonnegative reals
\item[(D)] $\{y \mid 0 < y \leq 1\}$
\end{list}

\bigskip\noindent\textbf{T24.}\ Which of the following is a reflection of the graph of $y = f(x)$ in the $y$-axis?
\begin{list}{}{\setlength{\leftmargin}{3em}\setlength{\labelwidth}{2.5em}\setlength{\itemsep}{0pt}\setlength{\topsep}{2pt}}
\item[(A)] $y = -f(x)$
\item[(B)] $y = f(-x)$
\item[(C)] $y = f(|x|)$
\item[(D)] $y = -f(-x)$
\end{list}

\bigskip\noindent\textbf{T25.}\ The smallest positive $x$ for which the function $f(x) = \sin\!\left(\dfrac{x}{4}\right) - 2$ is a maximum is
\begin{list}{}{\setlength{\leftmargin}{3em}\setlength{\labelwidth}{2.5em}\setlength{\itemsep}{0pt}\setlength{\topsep}{2pt}}
\item[(A)] $\dfrac{\pi}{2}$
\item[(B)] $\pi$
\item[(C)] $2\pi$
\item[(D)] $4\pi$
\end{list}

\bigskip\noindent\textbf{T26.}\ $\tan\!\left(\arccos\!\left(-\dfrac{\sqrt{3}}{2}\right)\right) =$
\begin{list}{}{\setlength{\leftmargin}{3em}\setlength{\labelwidth}{2.5em}\setlength{\itemsep}{0pt}\setlength{\topsep}{2pt}}
\item[(A)] $-1$
\item[(B)] $-\dfrac{\sqrt{3}}{3}$
\item[(C)] $\dfrac{\sqrt{3}}{3}$
\item[(D)] $1$
\end{list}

\bigskip\noindent\textbf{T27.}\ If $f^{-1}(x)$ is the inverse of $f(x) = 3e^{-x}$, then $f^{-1}(x) =$
\begin{list}{}{\setlength{\leftmargin}{3em}\setlength{\labelwidth}{2.5em}\setlength{\itemsep}{0pt}\setlength{\topsep}{2pt}}
\item[(A)] $\ln\!\left(\dfrac{3}{x}\right)$
\item[(B)] $\ln\!\left(\dfrac{x}{3}\right)$
\item[(C)] $\dfrac{1}{3}\ln x$
\item[(D)] $\sqrt{\ln x}$
\end{list}

\bigskip\noindent\textbf{T28.}\ Which of the following functions does not have an inverse function?
\begin{list}{}{\setlength{\leftmargin}{3em}\setlength{\labelwidth}{2.5em}\setlength{\itemsep}{0pt}\setlength{\topsep}{2pt}}
\item[(A)] $y = \cos x \ (0 \leq x \leq \pi)$
\item[(B)] $y = x^3 - 4$
\item[(C)] $y = \dfrac{x}{x^2+4}$
\item[(D)] $y = \ln(x+3)\ (\text{where } x > -3)$
\end{list}

\bigskip\noindent\textbf{T29.}\ Suppose that $f(x) = \ln x$ for all positive $x$ and $g(x) = 16 - x^2$ for all real $x$. The domain of $f(g(x))$ is
\begin{list}{}{\setlength{\leftmargin}{3em}\setlength{\labelwidth}{2.5em}\setlength{\itemsep}{0pt}\setlength{\topsep}{2pt}}
\item[(A)] $\{x \mid x \leq 4\}$
\item[(B)] $\{x \mid |x| > 4\}$
\item[(C)] $\{x \mid |x| < 4\}$
\item[(D)] $\{x \mid 0 < x < 4\}$
\end{list}

\bigskip\noindent\textbf{T30.}\ Suppose that $f(x) = \ln x$ for all positive $x$ and $g(x) = 16 - x^2$ for all real $x$. The range of $y = f(g(x))$ is
\begin{list}{}{\setlength{\leftmargin}{3em}\setlength{\labelwidth}{2.5em}\setlength{\itemsep}{0pt}\setlength{\topsep}{2pt}}
\item[(A)] $\{y \mid y > 0\}$
\item[(B)] $\{y \mid 0 < y \leq \ln 16\}$
\item[(C)] $\{y \mid y \leq \ln 16\}$
\item[(D)] $\{y \mid y < 0\}$
\end{list}

\bigskip\noindent\textbf{T31.}\ The curve defined parametrically by $x(t) = t^2 + 1$ and $y(t) = t^2 + 5$ is part of a(n)
\begin{list}{}{\setlength{\leftmargin}{3em}\setlength{\labelwidth}{2.5em}\setlength{\itemsep}{0pt}\setlength{\topsep}{2pt}}
\item[(A)] line
\item[(B)] circle
\item[(C)] parabola
\item[(D)] ellipse
\end{list}

\bigskip\noindent\textbf{T32.}\ Which equation includes the curve defined parametrically by $x(t) = \cos^2(t)$ and $y(t) = 3\sin(t)$?
\begin{list}{}{\setlength{\leftmargin}{3em}\setlength{\labelwidth}{2.5em}\setlength{\itemsep}{0pt}\setlength{\topsep}{2pt}}
\item[(A)] $x^2 + y^2 = 9$
\item[(B)] $9x^2 + y^2 = 9$
\item[(C)] $9x + y^2 = 9$
\item[(D)] $x + 9y^2 = 1$
\end{list}

\bigskip\noindent\textbf{T33.}\ Find the smallest value of $\theta$ in the interval $[0, 2\pi]$ for which the rose $r = 3\cos(4\theta)$ passes through the origin.
\begin{list}{}{\setlength{\leftmargin}{3em}\setlength{\labelwidth}{2.5em}\setlength{\itemsep}{0pt}\setlength{\topsep}{2pt}}
\item[(A)] $0$
\item[(B)] $\dfrac{\pi}{16}$
\item[(C)] $\dfrac{\pi}{8}$
\item[(D)] $\dfrac{\pi}{4}$
\end{list}

\bigskip\noindent\textbf{T34.}\ For what value of $\theta$ in the interval $[0,\pi]$ do the polar curves $r = 4$ and $r = 2 + 2\sqrt{2}\cos\theta$ intersect?
\begin{list}{}{\setlength{\leftmargin}{3em}\setlength{\labelwidth}{2.5em}\setlength{\itemsep}{0pt}\setlength{\topsep}{2pt}}
\item[(A)] $\dfrac{\pi}{6}$
\item[(B)] $\dfrac{\pi}{4}$
\item[(C)] $\dfrac{\pi}{3}$
\item[(D)] $\dfrac{2\pi}{3}$
\end{list}

\bigskip\noindent\textbf{TB1.}\ The graph of the function $f(x) = 3e^{\sin(x)} - 4$ crosses the $x$-axis once in the interval $[0,1]$. What is the $x$-coordinate of this $x$-intercept?
\begin{list}{}{\setlength{\leftmargin}{3em}\setlength{\labelwidth}{2.5em}\setlength{\itemsep}{0pt}\setlength{\topsep}{2pt}}
\item[(A)] $0.209$
\item[(B)] $0.292$
\item[(C)] $0.417$
\item[(D)] $0.552$
\end{list}

\bigskip\noindent\textbf{TB2.}\ Find the $x$-intercept of the graph of
\[ f(x) = \sqrt{\cos(x) + 1} + x^3 - 4e^{\sin(x)} + 7 \]
on the portion of the graph where $f(x)$ is decreasing.
\begin{list}{}{\setlength{\leftmargin}{3em}\setlength{\labelwidth}{2.5em}\setlength{\itemsep}{0pt}\setlength{\topsep}{2pt}}
\item[(A)] $-1.849$
\item[(B)] $-0.801$
\item[(C)] $0.984$
\item[(D)] $1.367$
\end{list}

\bigskip\noindent\textbf{TB3.}\ You are given the function $f(x) = e^{x^5 - 3x^3 + 1} - 5$ on the closed interval $[-2,2]$. Find all intervals where $f(x) < 0$.
\begin{list}{}{\setlength{\leftmargin}{3em}\setlength{\labelwidth}{2.5em}\setlength{\itemsep}{0pt}\setlength{\topsep}{2pt}}
\item[(A)] $[-2,-1.696)$ and $(-0.615,1.764)$
\item[(B)] $(-1.696,-0.615)$ only
\item[(C)] $[-2,-1.696)$ only
\item[(D)] $(-0.615,1.764)$ only
\end{list}

\bigskip\noindent\textbf{TB4.}\ You are given the function $f(x) = (6 - 2x - 2x^2)\cos(3x-2)$ on the closed interval $[-3,2]$. How many times does $f(x)$ cross the $x$-axis in the interval?
\begin{list}{}{\setlength{\leftmargin}{3em}\setlength{\labelwidth}{2.5em}\setlength{\itemsep}{0pt}\setlength{\topsep}{2pt}}
\item[(A)] four
\item[(B)] five
\item[(C)] six
\item[(D)] seven
\end{list}

\bigskip\noindent\textbf{TB5.}\ On the interval $[0,2\pi]$ there is one point on the curve $r = \theta - 3\cos\theta$ whose $x$-coordinate is $3$. Find the $y$-coordinate there.
\begin{list}{}{\setlength{\leftmargin}{3em}\setlength{\labelwidth}{2.5em}\setlength{\itemsep}{0pt}\setlength{\topsep}{2pt}}
\item[(A)] $-3.764$
\item[(B)] $-0.858$
\item[(C)] $1.979$
\item[(D)] $5.201$
\end{list}

\bigskip\noindent\textbf{TB6.}\ If $f(x) = 2 + e^x$ then the domain of $f^{-1}(x)$ is
\begin{list}{}{\setlength{\leftmargin}{3em}\setlength{\labelwidth}{2.5em}\setlength{\itemsep}{0pt}\setlength{\topsep}{2pt}}
\item[(A)] $(-\infty,\infty)$
\item[(B)] $(0,\infty)$
\item[(C)] $(1,\infty)$
\item[(D)] $(2,\infty)$
\end{list}

\end{document}
